\reportelement{Введение}

Цель работы --- подключить к приложению PostgreSQL, научиться хранить
данные вне контейнера и разобраться с сетями Docker.

Задачи:
\begin{itemize}
  \item создать именованный volume и запустить PostgreSQL с его использованием;
  \item подключить приложение к базе данных и проверить, что данные
        сохраняются после пересоздания контейнера;
  \item подключить файл конфигурации к контейнеру через bind mount;
  \item создать пользовательскую сеть и проверить сетевую изоляцию контейнеров.
\end{itemize}

Работа выполнялась на сервере linux-lab, Docker версии
29.1.3~\cite{docker-docs}.

\reportelement{Персональные параметры}

Персональные параметры приведены в таблице~\ref{tab:parameters}.

\begin{table}[H]
  \caption{Персональные параметры}\label{tab:parameters}
  \centering
  \begin{tabularx}{\textwidth}{|p{0.46\textwidth}|X|}
    \hline
    Параметр & Значение\\ \hline
    Номер ИСУ & 467697\\ \hline
    Персональный порт (assistant) & 8097\\ \hline
    Персональный порт (postgres) & 5527\\ \hline
    Персональный префикс & ter-stepanyan\\ \hline
    Секретный ключ & secret97ter-stepanyan\\ \hline
    Имя базы данных & ter-stepanyan97\\ \hline
  \end{tabularx}
\end{table}

Репозиторий с результатом (ветка \texttt{lab-3}): \url{\RepositoryURL}.

\chapter{Основная часть}

\section{Этап 1. Запуск контейнера с именованным volume}

Создаю именованный volume для БД и проверяю его создание на
рисунке~\ref{fig:1-1}.

\reportscreenshot{fig-1-1.png}{Создание и проверка тома}{fig:1-1}

Запуск PostgreSQL с моим volume и проверка статуса представлены на
рисунке~\ref{fig:1-2}.

\reportscreenshot{fig-1-2.png}{Запуск PostgreSQL}{fig:1-2}

Исходное приложение хранило заметки в списке в памяти процесса и не
использовало переменную \texttt{DATABASE\_URL}, поэтому после удаления
контейнера заметки пропадали. В \texttt{app/main.py} обработчики
\texttt{/notes} изменены так, чтобы заметки записывались в таблицу
PostgreSQL, а в \texttt{requirements.txt} добавлен драйвер \texttt{psycopg}.
Собираю образ приложения из \texttt{Dockerfile.optimized} на
рисунке~\ref{fig:1-3}.

\reportscreenshot{fig-1-3.png}{Сборка образа}{fig:1-3}

Запускаю ассистента и подключаю его к БД на рисунке~\ref{fig:1-4}.
Строка подключения передаётся через переменную \texttt{DATABASE\_URL}.

\reportscreenshot{fig-1-4.png}{Запуск ассистента и подключение к БД}{fig:1-4}

Проверяю работоспособность ассистента на рисунке~\ref{fig:1-5}: в ответе
\texttt{/notes} видна созданная заметка.

\reportscreenshot{fig-1-5.png}{Проверка работоспособности}{fig:1-5}

\section{Этап 2. Тест персистентности: пересоздание контейнера}

Удаляю контейнер с ассистентом и пересоздаю его на рисунке~\ref{fig:2-1}.

\reportscreenshot{fig-2-1.png}{Пересоздание ассистента}{fig:2-1}

Проверяю, что данные сохранились, и создаю артефакт на
рисунке~\ref{fig:2-2}. Заметка осталась, потому что она хранится в
PostgreSQL, а данные базы лежат в volume \texttt{ter-stepanyan-pgdata},
который не зависит от контейнера приложения.

\reportscreenshot{fig-2-2.png}{Артефакт персистентности}{fig:2-2}

\section{Этап 3. Bind mount для конфигурации}

Создаю папку \texttt{config}, в ней файл \texttt{app.conf} на
рисунке~\ref{fig:3-1}.

\reportscreenshot{fig-3-1.png}{Содержимое файла app.conf}{fig:3-1}

Запускаю контейнер с bind mount на рисунке~\ref{fig:3-2}: локальная папка
\texttt{config} монтируется внутрь контейнера по пути \texttt{/app/config}
в режиме только для чтения. Также проверяю конфигурацию внутри контейнера.

\reportscreenshot{fig-3-2.png}{Запуск контейнера с bind mount}{fig:3-2}

На хосте в файле \texttt{config/app.conf} значение
\texttt{maintenance\_mode} изменено на \texttt{true}, после чего контейнер
перезапущен командой \texttt{docker restart}. Проверка через
\texttt{docker exec} на рисунке~\ref{fig:3-3} показала, что внутри
контейнера файл содержит новое значение без пересборки образа.

\reportscreenshot{fig-3-3.png}{Проверка изменения значения}{fig:3-3}

\section{Этап 4. Пользовательские сети для изоляции}

Создаю пользовательскую сеть и удаляю старые контейнеры на
рисунке~\ref{fig:4-1}.

\reportscreenshot{fig-4-1.png}{Создание пользовательской сети}{fig:4-1}

На рисунке~\ref{fig:4-2} запускаю БД и ассистента в одной сети \\
\texttt{ter-stepanyan-backend}. Ассистент подключается к базе по имени
контейнера \texttt{ter-stepanyan-postgres} и порту 5432. База доступна
только внутри сети, снаружи к ней подключиться нельзя.

\reportscreenshot{fig-4-2.png}{Создание контейнеров для БД и ассистента}{fig:4-2}

Проверяю работоспособность, создаю новую заметку и получаю все заметки на
рисунке~\ref{fig:4-3}. Заметка с первого этапа тоже на месте.

\reportscreenshot{fig-4-3.png}{Проверка связности}{fig:4-3}

Для исследования сетевой изоляции запускаю временные контейнеры на
рисунке~\ref{fig:4-4}: первый вне сети, второй внутри сети. Контейнер вне
сети не находит \texttt{ter-stepanyan-postgres} по имени (ошибка
\texttt{bad address}), а контейнер внутри сети получает ответ на ping.

\reportscreenshot{fig-4-4.png}{Исследование сетевой изоляции}{fig:4-4}

\section{Этап 5. Финальный артефакт}

На рисунке~\ref{fig:5-1} представлена итоговая инфраструктура: контейнеры
\texttt{ter-stepanyan-postgres} и \texttt{ter-stepanyan-assistant}, volume
\texttt{ter-stepanyan-pgdata} и сеть \texttt{ter-stepanyan-backend}.

\reportscreenshot{fig-5-1.png}{Финальный артефакт}{fig:5-1}

\reportelement{Заключение}

В ходе работы приложение было подключено к PostgreSQL, а данные базы
вынесены в именованный volume. Проверено, что заметки сохраняются после
удаления и пересоздания контейнеров. С помощью bind mount файл
конфигурации подключён к контейнеру без пересборки образа. Пользовательская
сеть позволила контейнерам обращаться друг к другу по имени и изолировала
их от контейнеров вне сети.

\begin{thebibliography}{9}
  \bibitem{docker-docs}
  Docker Documentation: документация Docker.
  URL: \url{https://docs.docker.com/}
  (дата обращения: 26.10.2026).
\end{thebibliography}
